% Stand 25. Juni 2026
\subsection{Datenbereinigung}
\label{KapitelDatenanalyseDatenbereinigung}
Die pro Tag für den Großraum München verfügbaren Datenreihen
unterschiedlicher Sensoren und damit unterschiedlicher räumlicher
Position weisen eine stark heterogene Qualität auf. 
Aufgrund der pro Tag mit 80-100 Dateien 
(vgl.\,Kap.\,\ref{KapitelDatenanalyseDatenHerkunft}) eher geringen
Anzahl von Fein\-staub-Messungen ist es ratsam, möglichst
viele der Dateien nach einer eventuellen Datenbereinigung für 
die weiteren Analysen verwenden zu können und so wenig wie möglich 
Dateien zu verwerfen.

Die Messreihen der unterschiedlichen Sensoren lassen sich 
hinsichtlich Qualität grob in fünf Kategorien einteilen:
\begin{itemize}[topsep=2pt,parsep=1pt,itemsep=0.5pt]
\item[a)] die Messreihen weisen größere zeitliche Unterbrechungen auf;
\item[b)] die Messreihen weisen NA-Werte auf;
\item[c)] die Messreihen weisen im Vergleich zu Messreihen der übrigen Sensoren  
auffallend wenige verschiedene Messwerte auf;
\item[d)] die Messreihen weisen für den PM2.5- und/oder den PM10-Messwert 
\textit{aufeinanderfolgend} hohe Werte auf, die oberhalb der gesetzlichen Grenzwerte liegen;
\item[e)] die Messreihen sind von guter Qualität.
\end{itemize}

\begin{figure}[H] 
\centering
\includegraphics[scale = 0.51]{Datenaufbereitung_P1_Verlauf_mit_Toleranzband_Sensor1785} 
\caption{\small{Verlauf der $\mathrm{P}_{10}$-Messwerte eines Sensors am 2.\,April 2026. 
Etwa zwischen 10 und 14\,Uhr gab es eine Sequenz 
der Länge 92 von Messwerten zwischen $600-1100\mu g/m^3$ und damit deutlich oberhalb des Schwellwertes (\textcolor{red}{\rule[1mm]{3mm}{0.5mm}})}}
\label{Abb_Verlauf_mit_Toleranzband_Sensor1785_über_Schwellwert}
\end{figure} 

Abb.\,(\ref{Abb_Verlauf_mit_Toleranzband_Sensor1785_über_Schwellwert}) 
zeigt ein Beispiel der Kategorie d), 
Abb.\,(\ref{Abb_Verlauf_mit_Toleranzband_Sensor23496_mit_Unterbrechung}) 
ein Beispiel der Kategorie a).

\begin{figure}[H] 
\centering
\includegraphics[scale = 0.51]{Datenaufbereitung_P1_Verlauf_mit_Toleranzband_Sensor23496} 
\caption{\small{Verlauf der $\mathrm{P}_{10}$-Messwerte eines Sensors am 24.\,April 2026. 
Etwa von 9 bis 11\,Uhr gab es eine Unterbrechung der Mess-Aufzeichnung von insgesamt 103 Minuten Dauer}}
\label{Abb_Verlauf_mit_Toleranzband_Sensor23496_mit_Unterbrechung}
\end{figure} 

In beiden Abbildungen sind neben den Messwerten (\textcolor{brown}{$\bullet$}), 
noch der zentrierte gleitende Durchschnitt 
(\textcolor{NavyBlue}{\rule[1mm]{3mm}{0.5mm}}) 
und das Toleranzband
(\textcolor{YellowGreen}{\rule[-0.5mm]{3mm}{3mm}})  
eingezeichnet. Hierbei wird das Toleranzband durch das Intervall
$[\mathrm{gleitender\ Durchschnitt}\pm 2\cdot\mathrm{Standardabweichung\ der\ Messreihe}]$ definiert und dient wie auch der gleitende
Durchschnitt ausschließlich der visuellen Erkennung von
eventuellen Auffälligkeiten der Messereihe.\\

Mit Messreihen der Kategorien a) bis d) wurde folgendermaßen
verfahren:

\begin{itemize}[topsep=2pt,parsep=1pt,itemsep=0.5pt]
\item[a)] Messreihen mit Unterbrechungen der Mess-Aufzeichnung
von mehr als 180 Minruten Länge wurden für die weiteren Analysen
\textit{nicht} berücksichtigt.
\item[b)] Bei den Messdateien, in denen NA-Werte auftraten, 
handelt es sich um zeilenweise erratisch eingefügte
Messungen mit einer validen Zeitangabe aber fehlenden 
Messwerten. 
Nach der Entfernung dieser fehlerhaften Zeilen war 
die verbleibende Messreihe unauffällig und wurde für die
weiteren Auswertungen verwendet.  
\item[c)] Messreihen, die weniger als 20 verschiedene 
Messwerte pro Tag aufweisen, wurden \textit{nicht} weiter berücksichtigt. 
Vermutlich handelt es sich hier um einen Defekt oder eine 
Fehlkonfiguration des D/A-Wandlers des SDS011 Sensors.
\item[d)] Messreihen, bei denen mehr als 120 \textit{nacheinander folgende} 
Messwerte (Definition einer Sequenz der Länge 120) oberhalb 
eines Schwellwertes von
$100\mu g/m^3$ liegen, wurden \textit{nicht} weiter berücksichtigt. 
Andernfalls hätten solche Messungen die anschließend durchgeführten
Mittelwertbildungen und die darauf aufsetzenden Auswertungen stark 
verzerrt.
\end{itemize}

Für die getroffenen Festsetzungen gibt es keine wissenschaftliche 
Rechtfertigung; sie sind das Ergebnis einer Abwägung 
der antagonistischen Ziele einer möglichst großen Datenbasis 
einerseits und einer guten Datenqualität andererseits. \\

Durch diese Schritte der Datenbereinigung der täglichen 
Messdaten wurde die Anzahl der im weiteren Verlauf 
betrachteten, täglichen Messreihen um etwa 2-4 Stück, also grob
weniger als 5\% verkleinert. \\

Nach Abschluss der reproduzierbaren Datenbereinigung liegt
nunmehr ein konsistenter Datensatz vor, der im folgenden 
Kapitel zur Untersuchung der räumlichen Abhängig\-keits\-struk\-tur
verwendet wird.
  


